% !TeX root = ../informe.tex
%=======================================================================
%  config/preambulo.tex — paquetes, idioma y estilo visual
%
%  Nada de este archivo es específico del proyecto: solo define cómo se ve
%  el documento. Lo propio de la plantilla (cajas de instrucción,
%  marcadores de relleno, identificadores trazables, inclusión de
%  diagramas) vive en config/plantilla.tex.
%=======================================================================

%----------------------- Codificación e idioma -------------------------
\usepackage[T1]{fontenc}
\usepackage[utf8]{inputenc}
\usepackage[spanish,es-tabla,es-noquoting]{babel}   % es-tabla: «Tabla», no «Cuadro»
\usepackage{csquotes}                               % biblatex lo exige con babel
\usepackage{lmodern}
\usepackage{microtype}

% amsmath debe cargarse antes que cleveref, o cleveref aborta la compilación.
\usepackage{amsmath}
\usepackage{amssymb}

%------------------------------ Geometría ------------------------------
\usepackage[a4paper,top=2.6cm,bottom=2.6cm,left=2.8cm,right=2.5cm,
            headheight=16pt]{geometry}
\usepackage{setspace}
\onehalfspacing
\setlength{\parskip}{6pt}
\setlength{\parindent}{0pt}

%------------------------------- Tablas --------------------------------
\usepackage{booktabs}
\usepackage{array}
\usepackage{tabularx}
\usepackage{longtable}
\usepackage{xltabular}      % tabularx que además parte entre páginas
\usepackage{ragged2e}
\usepackage{multirow}
\newcolumntype{L}[1]{>{\RaggedRight\arraybackslash}p{#1}}
\newcolumntype{C}[1]{>{\Centering\arraybackslash}p{#1}}
\newcolumntype{Y}{>{\RaggedRight\arraybackslash}X}
\renewcommand{\arraystretch}{1.25}

%--------------------------- Gráficos y color --------------------------
\usepackage{graphicx}
\usepackage{adjustbox}
\graphicspath{{figuras/}}
\usepackage[table,dvipsnames]{xcolor}
\usepackage{tcolorbox}
\tcbuselibrary{skins,breakable}

% floatpag permite quitar el folio de la página de un float: \thispagestyle
% no vale ahí dentro, porque el estilo se fija al componer la página y no
% al encajonar el float. Lo necesitan los diagramas a página completa.
%
% ORDEN OBLIGATORIO: floatpag ANTES que float. Los dos redefinen \@float;
% al revés, floatpag pisa el parche de float y la opción [H] deja de
% existir («Unknown float option `H'» en cada tabla que la usa).
\usepackage{floatpag}
\usepackage{float}
\usepackage{needspace}
% Mantiene tablas y figuras dentro de la seccion donde se las presenta.
\usepackage[section]{placeins}

% Da prioridad a colocar los flotantes cerca de su texto de referencia.
\renewcommand{\topfraction}{.90}
\renewcommand{\bottomfraction}{.80}
\renewcommand{\textfraction}{.08}
\renewcommand{\floatpagefraction}{.75}
\setcounter{topnumber}{3}
\setcounter{bottomnumber}{2}
\setcounter{totalnumber}{5}
\usepackage{caption}
\usepackage{subcaption}
\captionsetup{font=small,labelfont=bf}
\usepackage{rotating}       % sidewaysfigure, para diagramas anchos
\usepackage{pifont}

%------------------------------ Listas ---------------------------------
\usepackage{enumitem}
\setlist{itemsep=2pt,parsep=0pt,topsep=4pt}

%--------------------------- Código fuente -----------------------------
\usepackage{listings}

\definecolor{codeFondo}{HTML}{F7F7F7}
\definecolor{codeComentario}{HTML}{6A737D}
\definecolor{codePalabra}{HTML}{0B5394}
\definecolor{codeCadena}{HTML}{9C27B0}
\definecolor{codeBorde}{HTML}{DDDDDD}

\lstdefinestyle{estilobase}{
    backgroundcolor = \color{codeFondo},
    basicstyle      = \ttfamily\footnotesize,
    commentstyle    = \color{codeComentario}\itshape,
    keywordstyle    = \color{codePalabra}\bfseries,
    stringstyle     = \color{codeCadena},
    numbers         = left,
    numberstyle     = \tiny\color{codeComentario},
    numbersep       = 8pt,
    frame           = single,
    rulecolor       = \color{codeBorde},
    breaklines      = true,
    breakatwhitespace = false,
    showstringspaces  = false,
    tabsize         = 2,
    captionpos      = b,
    xleftmargin     = 12pt,
    literate = % acentos y eñes dentro de los bloques de código
      {á}{{\'a}}1 {é}{{\'e}}1 {í}{{\'i}}1 {ó}{{\'o}}1 {ú}{{\'u}}1
      {Á}{{\'A}}1 {É}{{\'E}}1 {Í}{{\'I}}1 {Ó}{{\'O}}1 {Ú}{{\'U}}1
      {ñ}{{\~n}}1 {Ñ}{{\~N}}1 {ü}{{\"u}}1 {¿}{{?`}}1 {¡}{{!`}}1
}
\lstset{style=estilobase}

% Lenguajes que aparecen en este proyecto. Uso:
%   \begin{lstlisting}[language=Java,caption={...},label={lst:...}]
\lstdefinelanguage{yaml}{
    keywords        = {true,false,null,version,services,image,build,ports,
                       environment,volumes,networks,depends_on,profiles,name},
    keywordstyle    = \color{codePalabra}\bfseries,
    sensitive       = false,
    comment         = [l]{\#},
    morestring      = [b]",
    morestring      = [b]'
}
\lstdefinelanguage{json}{
    string          = [b]",
    morestring      = [b]',
    literate        = {:}{{{\color{codePalabra}:}}}1 {,}{{{\color{codePalabra},}}}1,
    sensitive       = false
}

%--------------------------- Estilo de página --------------------------
\definecolor{acento}{HTML}{1168BD}

\usepackage{fancyhdr}
\pagestyle{fancy}
\fancyhf{}
\fancyhead[L]{\small\nouppercase{\leftmark}}
\fancyhead[R]{\small\thepage}
\renewcommand{\headrulewidth}{0.4pt}

% La primera página de cada sección usa el estilo «plain»; se ajusta para
% que también muestre el número de página arriba a la derecha.
\fancypagestyle{plain}{%
    \fancyhf{}%
    \fancyhead[R]{\small\thepage}%
    \renewcommand{\headrulewidth}{0pt}%
}

% Este documento es «article»: el nivel superior es \section. Se le da el
% peso visual que en la plantilla original tenía \chapter.
\usepackage{titlesec}
\titleformat{\section}[hang]
  {\normalfont\Large\bfseries\color{acento}}
  {\thesection.}{10pt}{}
\titlespacing*{\section}{0pt}{18pt}{10pt}

% Los rótulos internos que no deben entrar al índice se escriben con
% \paragraph* (titlesec no respeta \setcounter{tocdepth} en los niveles
% que gestiona).

\setcounter{secnumdepth}{3}
\setcounter{tocdepth}{2}    % el índice llega hasta x.y; x.y.z no aporta

%----------------------- Referencias cruzadas --------------------------
\usepackage[hidelinks]{hyperref}
% El idioma va explícito: sin él, cleveref rotula «Section 6» en inglés
% aunque babel esté en español.
\usepackage[spanish,nameinlink,capitalise]{cleveref}
% cleveref en español rotula las tablas como «Cuadro», pero babel con la
% opción es-tabla titula los pies como «Tabla». Sin esto, la prosa diría
% «Cuadro 3» y el pie de la misma tabla, «Tabla 3». Van capitalizadas las
% dos para no anular la opción «capitalise».
\crefname{table}{Tabla}{Tablas}
\Crefname{table}{Tabla}{Tablas}

%------------------------- Bibliografía --------------------------------
% Estilo IEEE, habitual en ingeniería. Requiere biber, que latexmk ya
% ejecuta solo.
\usepackage[backend=biber,style=ieee,sorting=nyt]{biblatex}
\addbibresource{referencias.bib}
